\section{Foreword}

\textit{If a principle may not be tested, then it may not truly be believed.}

Let's begin with a quote from John Evjen that appeared in a 1930s book on Lutheranism, "It was my good fortune to come to Augsburg College and Theological Seminary, Minneapolis, headed by the greatest theologian that America has ever had, Professor Georg Sverdrup...After completing seven years in this institution..."\footnote{John O. Evjen, ``What Is Lutheranism?'' in \textit{What Is Lutheranism? A Symposium in Interpretation}, ed. Vergilius Ferm (New York: The Macmillan Company, 1930), 9.}

Never forget that statement as you read the remainder of this volume. Evjen adored and respected Sverdrup. This was true in 1910 just as is was in 1930.   

Evjen was an engineer at heart. I see many of my engineering instincts in his writings. He was fastidious to a fault. He was a man who latched onto a single idea and worked on it for years. The source texts were consulted at length; dog-eared pages, notes in the margins, piles of notes. The idea is placed under a microscope, twisted, pulled, broken, and reassembled until there is clarity.  

Evjen dedicated a significant portion of his life defending Sverdrup's work. The contents of this book document parts of his work from 1910 through 1918. I would argue that significant work was done before 1910, since Evjen's 1910 polemic understood the Sverdrup P1 controversy from the very first line as a quote from Bjørnstjerne Bjørnson. It a prophetic poem for both Sverdrup and Evjen. 

\begin{quote}
When all that you meant was met with mockery,\\
and all that you wrought was only despised,\\
when your whole intent, however good it was,\\
is dismissed as nonsense by a self-satisfied fool,\\
then it must be practiced,\\
then it shall be tried,\\
your love.\\
\end{quote}

Sverdrup and, later, Evjen struggled to define the boundaries between faithful adherence to the Word, dogmatics, ecclesiology, and polity. All of this is done at the congregation level, which P1 defines as "the right form of the Kingdom of God on earth."

With that said, why is Evjen such a controversial figure in AFLC history?

Evjen was a systematist attempting to carefully define his terms and then place the ideas within a greater structure, document the consequences, and place clear warning labels on the load-bearing beams. Make no mistake: you and I as faithful members of the AFLC are standing on those same load-bearing beams today.

Evjen asks the same question we should ask ourselves today: can we identify the beams, and do we understand how much weight they carry? What happens if one fails or if we start to move the load-bearing walls? 

As a test case, consider Gene Edward Veith's work on vocation. In practice, we often think and act in what Luther calls the mask of God. Suppose we amend the Fundamental Principles so that they explicitly call the individual to vocation. This is a dangerous ecclesiology plus polity question unless we fully understand the structure and order preserved by the Principles. 

Vocation is also ecclesial, not merely individual, a point developed later through the sailor's model. This is the clarity passed down from Sverdrup and Oftedal. P1 demands nothing less than all God's gifts by all hands within the congregation. This is the visible congregation manifesting spirit and life. The struggle for a clear definition of congregation is uncomfortable gift that John Evjen has given the AFLC. With an engineer's mind he stress-tests the Fundamental Principles starting with the first:
 
\begin{quote}
According to the Word of God, the congregation is the right form of the Kingdom of God on earth. 
\end{quote}

Before we continue, we must understand that Evjen is a flawed witness to the AFLC principles. In this book you will see his academic tenacity, his imbalance that many found disturbing, his unwillingness to yield, and language that is downright offensive to a modern audience. He is no hero, any more than Sverdrup or Oftedal were heroes. Yet, these flawed men speak to us today. We stand in wonder at the boundaries they established for our congregations.

But like a newly discovered fence in the woods, we may forget why it is there. We should therefore be very careful and not remove Chesterton's fence. We do not know if it is keeping something out, purposefully constraining those within its boundaries, or serving as a visible boundary that gives shape and meaning to something precious.  

Yet, P1 is too strong a claim to survive on sentiment. Without a living apologetic, it will either be rejected by wider Lutheran instincts or reduced within the AFLC to the vague language of “spirit and life.” Without the inherited obligation, P1 becomes the very thing Sverdrup and Evjen warn us against over and over again. 

This work serves to breathe life into the Principles. P1 is not presented as a clean description but as a living idea held upright in tension. 

Evjen is necessary to reveal the foundation of the AFLC.  
